\chapter{1977年内蒙古卷（理）}

\section{解答题}

\begin{problem}
\begin{enumerate}
\item 化简：\(\sqrt{4-12a+9a^2}\).

\item 计算：\(\dfrac{\left(-27\right)^3\sqrt{\left(\dfrac{1}{16}\right)^{-\frac{1}{2}}}}{\left(-9\right)^0\left(\dfrac{3}{2}\right)^7\left(\dfrac{1}{2}\right)^{-5}}\).

\item 已知 \(\lg2=0.3010\)，\(f\left(x\right)=\log_{\sqrt{10}}x\)，求 \(f\left(5\right)\).

\item 求函数 \(y=\dfrac{\sqrt{1-\frac{2}{3}x}}{\sqrt[3]{3x-1}}\) 的定义域.

\item 已知 \(\cos\theta=0.6\)，\(270^\circ<\theta<360^\circ\)，求 \(\left(\sin\dfrac{\theta}{2}-\cos\dfrac{\theta}{2}\right)^2\) 的值.

\item 把点 \(P\) 的极坐标 \(\left(-4,-\dfrac{4\pi}{3}\right)\) 化成直角坐标.
\end{enumerate}
\end{problem}

\begin{solution}
\begin{enumerate}
\item 先配成完全平方式，得 \(4-12a+9a^2=\left(2-3a\right)^2\). 因为平方根表示算术平方根，所以 \(\sqrt{\left(2-3a\right)^2}=\left|2-3a\right|\). 当 \(a\leq\frac{2}{3}\) 时，\(2-3a\geq0\)，当 \(a>\frac{2}{3}\) 时，\(2-3a<0\)，故
\[
\sqrt{4-12a+9a^2}=\left|2-3a\right|=
\begin{cases}
2-3a,&a\leq\frac{2}{3},\\
3a-2,&a>\frac{2}{3}.
\end{cases}
\]

\item 因为 \((-27)^3=-3^9\)，\(\left(\dfrac{1}{16}\right)^{-\frac12}=4\)，所以
\[
\begin{gathered}
\frac{(-27)^3\sqrt{\left(\frac{1}{16}\right)^{-\frac12}}}{(-9)^0\left(\frac32\right)^7\left(\frac12\right)^{-5}}
=\frac{-3^9\cdot2}{1\cdot\frac{3^7}{2^7}\cdot2^5}
=\frac{-3^9\cdot2}{3^7\cdot2^{-2}}\\
=-3^2\cdot2^3=-72.
\end{gathered}
\]

\item 由换底公式，\(f\left(5\right)=\dfrac{\lg5}{\lg\sqrt{10}}\). 又因为 \(\lg5=\lg\dfrac{10}{2}=\lg10-\lg2=1-0.3010=0.6990\)，且 \(\lg\sqrt{10}=\dfrac12\lg10=\dfrac12\)，所以
\[
f\left(5\right)=\frac{\lg5}{\lg\sqrt{10}}=\frac{1-0.3010}{\frac12}=1.398
\]

\item 分式有意义，必须同时满足根号内非负和分母不为零，即
\[
\begin{cases}
1-\dfrac23x\geq0,\\
3x-1\neq0.
\end{cases}
\]
第一式给出 \(x\leq\dfrac32\)，第二式给出 \(x\neq\dfrac13\). 因此定义域为
\[
\left(-\infty,\frac13\right)\cup\left(\frac13,\frac32\right]
\]

\item 因为 \(270^\circ<\theta<360^\circ\)，所以 \(\sin\theta<0\). 由 \(\sin^2\theta+\cos^2\theta=1\)，得 \(\sin\theta=-\sqrt{1-0.6^2}=-0.8\). 因此
\[
\begin{gathered}
\left(\sin\frac\theta2-\cos\frac\theta2\right)^2
=\sin^2\frac\theta2+\cos^2\frac\theta2-2\sin\frac\theta2\cos\frac\theta2\\
=1-\sin\theta=1+0.8=1.8.
\end{gathered}
\]

\item 设极坐标为 \((\rho,\varphi)\)，则直角坐标满足 \(x=\rho\cos\varphi\)，\(y=\rho\sin\varphi\). 对于 \(\rho=-4\)，\(\varphi=-\dfrac{4\pi}{3}\)，有
\[
x=-4\cos\left(-\frac{4\pi}{3}\right)=-4\cos\left(\pi+\frac\pi3\right)=-4\left(-\frac12\right)=2
\]
又因为 \(\sin\left(-\dfrac{4\pi}{3}\right)=-\sin\dfrac{4\pi}{3}=\dfrac{\sqrt3}{2}\)，所以
\[
y=-4\sin\left(-\frac{4\pi}{3}\right)=-4\cdot\frac{\sqrt3}{2}=-2\sqrt3
\]
故点 \(P\) 的直角坐标为 \(\left(2,-2\sqrt3\right)\).
\end{enumerate}
\end{solution}

\begin{problem}
\begin{enumerate}
\item 已知梯形 \(ABCD\)，\(AD\parallel BC\)，对角线 \(AC\) 和 \(BD\) 相交于点 \(E\)，过点 \(E\) 作平行于底的直线，交 \(AB\) 于点 \(M\)，交 \(CD\) 于点 \(N\)，求证：\(ME=EN\).

\item 圆台上底直径为 \(12\mathrm{~cm}\)，下底直径为 \(24\mathrm{~cm}\)，高为 \(8\mathrm{~cm}\)，求圆台的全面积.（\(\pi=3.14\)，精确到 \(0.1\mathrm{~cm}^2\)）
\end{enumerate}
\end{problem}

\begin{solution}
\begin{enumerate}
\item 由题意，\(AD\parallel MN\parallel BC\). 在两条截线上应用平行线分线段成比例定理，得
\[
\frac{AM}{AB}=\frac{DN}{DC}
\]

\solutionfigure{\bitmapfigure[max width=4.9cm]{img/1977/inner_mongolia_science/q2_solution_fig1.png}}

在 \(\triangle ABC\) 中，因 \(ME\parallel BC\)，所以 \(\triangle AME\sim\triangle ABC\)，从而 \(\dfrac{ME}{BC}=\dfrac{AM}{AB}\). 在 \(\triangle DBC\) 中，因 \(EN\parallel BC\)，所以 \(\triangle DEN\sim\triangle DBC\)，从而 \(\dfrac{EN}{BC}=\dfrac{DN}{DC}\). 结合前面的比例式，得到 \(\dfrac{ME}{BC}=\dfrac{EN}{BC}\)，故 \(ME=EN\).

\item 上、下底半径分别为 \(r_1=\dfrac12\times12\mathrm{~cm}=6\mathrm{~cm}\) 和 \(r_2=\dfrac12\times24\mathrm{~cm}=12\mathrm{~cm}\)，高为 \(h=8\mathrm{~cm}\). 圆台的母线长为
\[
L=\sqrt{h^2+\left(r_2-r_1\right)^2}=\sqrt{8^2+\left(12-6\right)^2}=10\mathrm{~cm}
\]
圆台的全面积等于两个底面面积与侧面积之和，因此
\[
S=\pi\left[r_1^2+r_2^2+\left(r_1+r_2\right)L\right]
\approx3.14\left[6^2+12^2+\left(6+12\right)\times10\right]
=3.14\times360=1130.4\mathrm{~cm}^2
\]
所以圆台的全面积约为 \(1130.4\mathrm{~cm}^2\).
\end{enumerate}
\end{solution}

\begin{problem}
红旗拖拉机厂今年元月份生产出一批甲、乙两种型号的拖拉机，其中生产乙型 \(16\text{ 台}\). 从二月份起，甲型每月增产 \(10\text{ 台}\)，而乙型按每月相同增长率逐步递增. 又知二月份甲、乙两型拖拉机产量之比是 \(3:2\)，三月份甲、乙两型拖拉机产量之和是 \(65\text{ 台}\). 求乙型拖拉机每月的增长率及甲型拖拉机元月份的产量.
\end{problem}

\begin{solution}
设乙型拖拉机每月的增长率为 \(x\)，甲型拖拉机元月份的产量为 \(y\text{ 台}\). 因为乙型二月份和三月份的产量分别为 \(16\left(1+x\right)\) 和 \(16\left(1+x\right)^2\)，甲型二月份和三月份的产量分别为 \(y+10\) 和 \(y+20\)，所以依题意得
\[
\begin{cases}
\dfrac{y+10}{16\left(1+x\right)}=\dfrac32,\\
y+20+16\left(1+x\right)^2=65.
\end{cases}
\]
第一式化为 \(y+10=24\left(1+x\right)\)，即 \(y=24x+14\). 代入第二式，得
\[
16x^2+56x-15=0
\]
因式分解得 \(\left(4x-1\right)\left(4x+15\right)=0\).
所以 \(x=\dfrac14\) 或 \(x=-\dfrac{15}{4}\). 由于产量按增长率递增，必须有 \(1+x>0\)，且元月份产量 \(y\) 应为正数，故舍去 \(x=-\dfrac{15}{4}\)，取 \(x=\dfrac14=25\%\). 此时
\[
y=24\times\frac14+14=20
\]
验算：二月份两型产量为 \(30\text{ 台}\) 和 \(20\text{ 台}\)，比为 \(3:2\)；三月份两型产量为 \(40\text{ 台}\) 和 \(25\text{ 台}\)，和为 \(65\text{ 台}\)，满足题意. 因此乙型拖拉机每月的增长率为 \(25\%\)，甲型拖拉机元月份的产量为 \(20\text{ 台}\).
\end{solution}

\begin{problem}
已知直线 \(3x+y=5\) 及 \(2x-3y+4=0\) 交于点 \(A\). 求：
\begin{enumerate}
\item 过点 \(A\) 且与 \(y\) 轴垂直的直线方程；
\item 过点 \(A\) 且与圆 \(x^2+y^2=1\) 相切的直线方程.
\end{enumerate}
\end{problem}

\begin{solution}
点 \(A\) 同时在两条已知直线上，所以其坐标满足
\[
\begin{cases}
3x+y=5,\\
2x-3y+4=0.
\end{cases}
\]
由第一式得 \(y=5-3x\)，代入第二式得 \(2x-3\left(5-3x\right)+4=0\)，即 \(11x-11=0\)，所以 \(x=1\)，\(y=2\). 故 \(A\left(1,2\right)\).

\begin{enumerate}
\item 与 \(y\) 轴垂直的直线是水平直线，过 \(A\left(1,2\right)\) 的方程为 \(y=2\).

\item 设切点为 \(P\left(u,v\right)\). 因为 \(P\) 在圆上，所以 \(u^2+v^2=1\). 半径 \(OP\) 垂直于切线 \(AP\)，故向量 \(\overrightarrow{OP}=\left(u,v\right)\) 与向量 \(\overrightarrow{PA}=\left(1-u,2-v\right)\) 的内积为零，即
\[
u(1-u)+v(2-v)=0
\]
结合 \(u^2+v^2=1\)，化简得 \(u+2v=1\). 于是 \(u=1-2v\)，代入圆方程得
\(\left(1-2v\right)^2+v^2=1\)，\(5v^2-4v=0\)，\(v\left(5v-4\right)=0\).
因此切点有两个：\(P_1\left(1,0\right)\) 和 \(P_2\left(-\dfrac35,\dfrac45\right)\).

圆 \(x^2+y^2=1\) 在点 \(P\left(u,v\right)\) 处的切线方程为 \(ux+vy=1\). 对于 \(P_1\)，切线为 \(x=1\)；对于 \(P_2\)，切线为 \(-\dfrac35x+\dfrac45y=1\)，即 \(3x-4y+5=0\). 所求直线为 \(x=1\) 和 \(3x-4y+5=0\).
\end{enumerate}
\end{solution}

\begin{problem}
甲、乙两船，甲船在某岛 \(B\) 的正南方向 \(A\) 处，\(AB=10\text{ 海里}\)，甲船自 \(A\) 处以 \(4\text{ 海里/小时}\) 的速度向正北方向航行；同时，乙船以 \(6\text{ 海里/小时}\) 的速度自岛 \(B\) 出发，向岛的北 \(60^\circ\) 西方向驶去. 问几分钟后两船相距最近？（精确到 \(1\text{ 分}\)）
\end{problem}

\begin{solution}
设经过 \(t\) 小时后，甲船航行到 \(C\) 点，乙船航行到 \(D\) 点，并取 \(B\) 为原点、正北方向为 \(y\) 轴正方向. 则
\(C\left(0,4t-10\right)\)，\(D\left(-6t\sin60^\circ,6t\cos60^\circ\right)=\left(-3\sqrt3t,3t\right)\).

\solutionfigure{\bitmapfigure[max width=3.8cm]{img/1977/inner_mongolia_science/q5_solution_fig1.png}}

无论甲船是否已经越过 \(B\)，两船距离的平方都可以由坐标直接求得：
\[
\begin{gathered}
DC^2=\left(3\sqrt3t\right)^2+\left(3t-(4t-10)\right)^2
=27t^2+(10-t)^2\\
=28t^2-20t+100
=28\left(t-\frac5{14}\right)^2+\frac{675}{7}.
\end{gathered}
\]
因为 \(t\geq0\)，且 \(\dfrac5{14}>0\)，所以当 \(t=\dfrac5{14}\) 小时，\(DC^2\) 取得最小值，从而 \(DC\) 取得最小值. 换算为分钟，\(\dfrac5{14}\times60=\dfrac{150}{7}\approx21.4\) 分钟，按题意精确到 \(1\) 分为 \(21\) 分钟.
\end{solution}

\begin{problem}
已知在 \(\triangle ABC\) 中，\(a\)，\(b\)，\(c\) 是三个内角 \(A\)，\(B\)，\(C\) 的对边，\(\lg\sin A\)，\(\lg\sin B\)，\(\lg\sin C\) 成等差数列.
\begin{enumerate}
\item 求证：\(\dfrac{\sin^2 A}{\sin^2 B}=\dfrac{a}{c}\)；

\item 又若方程 \(cx^2+2cx+a=0\) 有相等两实根，求证：\(\sin A=\sin B=\sin C\).
\end{enumerate}
\end{problem}

\begin{solution}
\begin{enumerate}
\item 三角形内角均在 \(\left(0,\pi\right)\) 内，所以各正弦值均为正数，取对数不会改变等式的合法性. 由三项成等差数列，得
\[
2\lg\sin B=\lg\sin A+\lg\sin C
\]
从而 \(\sin^2B=\sin A\sin C\). 由正弦定理，\(\dfrac{a}{\sin A}=\dfrac{c}{\sin C}\)，即 \(\sin C=\dfrac{c}{a}\sin A\). 代入前式，得
\[
\sin^2B=\frac{c}{a}\sin^2A
\]
所以 \(\dfrac{\sin^2A}{\sin^2B}=\dfrac{a}{c}\).

\item 方程有相等两实根，所以判别式为零：
\[
\Delta=\left(2c\right)^2-4ca=4c\left(c-a\right)=0
\]
由于 \(c\) 是三角形边长，\(c>0\)，故 \(c=a\). 由正弦定理，\(a=c\) 得 \(\sin A=\sin C\). 再由前面得到的关系 \(\sin^2B=\sin A\sin C\)，可得 \(\sin^2B=\sin^2A\). 因为 \(A\)、\(B\) 是三角形内角，正弦值均为正数，所以 \(\sin B=\sin A\). 因此 \(\sin A=\sin B=\sin C\).
\end{enumerate}
\end{solution}

\section{附加题}

\begin{problem}
在 \(\triangle ABC\) 中，求证 \(\sin^2 A+\sin^2 B+\cos^2 C+2\sin A\sin B\cos\left(A+B\right)=1\).
\end{problem}

\begin{solution}
因为 \(A+B+C=\pi\)，所以 \(\cos^2C=\cos^2\left(A+B\right)\). 因此左边为
\[
\begin{gathered}
\sin^2A+\sin^2B+\cos^2\left(A+B\right)+2\sin A\sin B\cos\left(A+B\right)\\
=\sin^2A+\sin^2B+\left(\cos A\cos B-\sin A\sin B\right)^2\\
\quad+2\sin A\sin B\left(\cos A\cos B-\sin A\sin B\right)\\
=\sin^2A+\sin^2B+\cos^2A\cos^2B-\sin^2A\sin^2B\\
=\sin^2A\left(1-\sin^2B\right)+\sin^2B+\cos^2A\cos^2B\\
=\sin^2A\cos^2B+\sin^2B+\cos^2A\cos^2B\\
=\cos^2B\left(\sin^2A+\cos^2A\right)+\sin^2B=1.
\end{gathered}
\]
故等式成立.
\end{solution}

\begin{problem}
已知椭圆的中心在原点，长轴和短轴都在坐标轴上，离心率 \(e=0.8\)，一条准线方程为 \(y=-\dfrac{25}{4}\)，求内接于椭圆的最大矩形面积.
\end{problem}

\begin{solution}
准线 \(y=-\dfrac{25}{4}\) 平行于 \(x\) 轴，故椭圆的长轴在 \(y\) 轴上、短轴在 \(x\) 轴上. 设半长轴为 \(a\)，半短轴为 \(b\)，半焦距为 \(c\)，则椭圆方程为
\(\frac{x^2}{b^2}+\frac{y^2}{a^2}=1\)，\(a>b>0\).
由离心率和准线方程，得
\(\frac{c}{a}=0.8=\frac45\)，\(\frac{a^2}{c}=\frac{25}{4}\).
于是 \(c=\dfrac45a\)，代入第二式得 \(\dfrac54a=\dfrac{25}{4}\)，所以 \(a=5\)，\(c=4\). 再由 \(b^2=a^2-c^2\)，得 \(b=\sqrt{25-16}=3\)，椭圆方程为
\[
\frac{x^2}{3^2}+\frac{y^2}{5^2}=1
\]

下面求内接矩形的最大面积. 椭圆的方程写成 \(Q\left(x,y\right)=\dfrac{x^2}{b^2}+\dfrac{y^2}{a^2}=1\). 若内接矩形中心为 \(M\)，取其一条对角线方向向量为 \(w\)，则两个对角顶点为 \(M+w\) 和 \(M-w\)，它们都在椭圆上. 将这两个点分别代入 \(Q=1\) 并相减，得
\[
\frac{M_xw_x}{b^2}+\frac{M_yw_y}{a^2}=0
\]
两条对角线方向向量不共线，因此对两条对角线同时成立时只能有 \(M_x=M_y=0\)，即矩形中心为椭圆中心.

设矩形的两个半边向量为 \(\left(u_1,u_2\right)\) 和 \(\left(v_1,v_2\right)\). 矩形的四个顶点为 \(\pm\left(u_1,u_2\right)\pm\left(v_1,v_2\right)\)，且两半边向量垂直，所以 \(u_1v_1+u_2v_2=0\). 又因为顶点 \(\left(u_1,u_2\right)+\left(v_1,v_2\right)\) 和 \(\left(u_1,u_2\right)-\left(v_1,v_2\right)\) 都在椭圆上，代入椭圆方程并相减，得
\[
\frac{u_1v_1}{b^2}+\frac{u_2v_2}{a^2}=0
\]
结合 \(u_1v_1+u_2v_2=0\) 及 \(a\neq b\)，可得 \(u_1v_1=0\)，\(u_2v_2=0\). 因此矩形的边平行于坐标轴. 设其半长、半宽分别为 \(p\)、\(q\)，则
\[
\frac{p^2}{b^2}+\frac{q^2}{a^2}=1
\]
且 \(S=4pq\).
由均值不等式，
\[
1=\frac{p^2}{b^2}+\frac{q^2}{a^2}\geq2\frac{pq}{ab}
\]
所以 \(S=4pq\leq2ab\). 当 \(\dfrac{p}{b}=\dfrac{q}{a}=\dfrac1{\sqrt2}\) 时取等号，此时矩形的两边分别为 \(2p=3\sqrt2\) 和 \(2q=5\sqrt2\)，故最大面积为
\[
S_{\max}=2ab=2\times3\times5=30
\]
\end{solution}
